% ============================================================
%  SECTION 4: DATA
% ============================================================

\section{Data}
\label{sec:data}

\subsection{Chetty et al.\ (2020) Panel}

We use the college-level panel constructed by
\textcite{chetty2020mobility}, which links tax records and college
enrollment data to measure the parental income distribution of students
at each institution. The panel covers 2,202 colleges across 52 states
and territories for birth cohorts 1980--1991 (approximately
corresponding to college entry years 1998--2009). Students are assigned
to the institution attended at ages 19--22.

The primary analysis table (MRC Table~3) provides college-by-cohort
observations with measures of parental income rank and quintile shares.
We supplement this with institutional characteristics from MRC Table~10
(SAT scores, expenditures, enrollment, Barron's selectivity, HBCU and
flagship status, graduation rates). These characteristics capture the
dimensions of institutional quality that prior work has linked to
earnings premia \parencite{brewer1999elite, dale2002selective}.

\subsection{Sample Construction}

We define treated institutions as Tennessee public colleges. Tennessee
has 51 colleges in the data, of which 26 are public---19 two-year
institutions and 7 four-year institutions. Multi-campus systems appear
as distinct identifiers (56 per cohort), but we identify 26 unique
public colleges.

The control group comprises public colleges in states that did not adopt
broad-based merit scholarship programs during the analysis period. We
drop all 25 states classified as merit-program adopters by
\textcite{sjoquist2015state}. The remaining 26 control states contribute
6,289 college-cohort observations. Of the 312 possible Tennessee
public-college-by-cohort observations ($26 \times 12$), 291 are
non-missing on both primary outcomes.

We exclude 24 observations with missing state or institution type.

\subsection{Outcome Variables}

Our two primary outcomes are:

\begin{itemize}
\item \textbf{Bottom-quintile share} ($Q_1$): the fraction of students
  at college $c$ in cohort $t$ whose parents are in the bottom 20\% of
  the national income distribution. Baseline mean for Tennessee public
  colleges: 0.138.

\item \textbf{Mean parental income rank} ($\overline{r}$): the average
  percentile rank of parental household income for students at college
  $c$ in cohort $t$. Baseline mean for Tennessee publics: 0.540.
\end{itemize}

A clarification on the unit of observation: each outcome is measured at
the college-by-birth-cohort level, not for the entire enrolled student
body at a point in time. For example, $Q_1$ for the University of
Tennessee in cohort 1986 measures the fraction of all 1986-born students
who attended UT between ages 19 and 22 whose parents were in the bottom
quintile. The Fall 2006 student body at UT, by contrast, would include
birth cohorts 1984--1987 simultaneously---mixing untreated and treated
cohorts. The cohort-specific measure provides sharper identification
because it isolates students who faced HOPE-era incentives from those
who did not. By the end of the post-treatment window ($k = +5$, cohort
1991), nearly all cohorts enrolled at a given institution are treated, so
cohort-level and student-body-level effects converge.

For the quintile decomposition, we additionally use the share of
students from each income quintile ($Q_1$ through $Q_5$).

For the level decomposition that addresses ecological inference
concerns, we construct count outcomes by multiplying quintile shares by
cohort size. For example, Q1 count $= Q_1 \times N_{ct}$, where
$N_{ct}$ is total enrollment at college $c$ in cohort $t$.

\subsection{Summary Statistics}

Table~\ref{tab:summary_stats} reports summary statistics for the primary
outcomes and quintile shares, separately for Tennessee public colleges
and control-state public colleges. Tennessee publics have a somewhat
higher baseline share of bottom-quintile students (14.0\% vs.\ 13.6\%)
and lower mean parental income rank (0.524 vs.\ 0.541). The quintile
distribution at Tennessee publics is shifted leftward: Q1 and Q2 shares
are higher, Q5 share is lower. Table~\ref{tab:did_main} reports pre-
and post-treatment means. After HOPE implementation, bottom-quintile
share at Tennessee publics rises to 15.9\%, while the control group
increases modestly to 12.9\%. The raw difference-in-differences is 1.5
percentage points.

\begin{table}[H]
\centering
\begin{threeparttable}
\caption{Summary Statistics}
\label{tab:summary_stats}
\begin{tabular}{lcccc}
\toprule
 & \multicolumn{2}{c}{TN Public} & \multicolumn{2}{c}{Control Public} \\
\cmidrule(lr){2-3} \cmidrule(lr){4-5}
 & Mean & SD & Mean & SD \\
\midrule
\multicolumn{5}{l}{\textit{Primary Outcomes}} \\
\quad Bottom-quintile share ($Q_1$) & 0.140 & 0.046 & 0.136 & 0.083 \\
\quad Mean parental income rank & 0.524 & 0.067 & 0.541 & 0.103 \\
\addlinespace
\multicolumn{5}{l}{\textit{Quintile Shares}} \\
\quad Q1 (bottom) & 0.140 & & 0.136 & \\
\quad Q2 & 0.195 & & 0.182 & \\
\quad Q3 & 0.239 & & 0.227 & \\
\quad Q4 & 0.256 & & 0.252 & \\
\quad Q5 (top) & 0.170 & & 0.203 & \\
\addlinespace
Observations & \multicolumn{2}{c}{137} & \multicolumn{2}{c}{2,850} \\
Colleges & \multicolumn{2}{c}{26} & \multicolumn{2}{c}{$\sim$500} \\
\bottomrule
\end{tabular}
\begin{tablenotes}[flushleft]
\footnotesize
\item \textit{Notes:} Pre-treatment means (birth cohorts 1980--1985)
  for Tennessee public colleges (treated) and public colleges in 26
  non-merit-program states (control). Quintile shares refer to the
  fraction of students whose parents are in the indicated quintile of
  the national income distribution. Observations are at the
  college-by-cohort level. Data: \textcite{chetty2020mobility}.
\end{tablenotes}
\end{threeparttable}
\end{table}
